\par\Needspace{6\baselineskip}\medskip
\begin{center}\hypertarget{native-179}{}\hypertarget{native-179-article132}{}\pdfbookmark[2]{132.}{article-132}132.\end{center}
For brevity, let us put
\[
\begin{gathered}
(0.1.2)=a,\qquad (0.\mathrm{I}.2)D'=-b,\\
(0.\mathrm{O}.2)D=c,\qquad (0.\mathrm{II}.2)D''=d,
\end{gathered}
\]
so that equation 10 of Article 114 becomes
\[
a\delta'=b+c\cdot\frac{n}{n'}+d\cdot\frac{n''}{n'}\mathpunct{\raisebox{.45ex}{,}}
\]
where the coefficients $c$ and $d$ will be of the first order and the difference $c-d$ is readily shown
to be of the second order. It follows that the value of $\dfrac{cn+dn''}{n+n''}$ obtained from the approximate assumption
$n:n''=\theta:\theta''$ has an error of only the fourth order, or even only the fifth whenever the middle observation is
separated from the two outer ones by equal intervals. This error is
\[
=\frac{c\theta+d\theta''}{\theta+\theta''\mathpunct{\mkern-8mu\raisebox{-.65ex}{.}\mkern8mu}}
-\frac{cn+dn''}{n+n''}
=\frac{\theta\theta''(d-c)(\eta''-\eta)}{(\theta+\theta'')(\eta''\theta+\eta\theta'')}\mathpunct{\raisebox{.45ex}{,}}
\]
where the denominator is of the second order, and one factor in the numerator, $\theta\theta''(d-c)$, is of the
fourth order, while the other, $\eta''-\eta$, is of the second order or, in the special case just mentioned,
of the third. Writing the equation in the form
\[
a\delta'=b+\frac{cn+dn''}{n+n''}\cdot\frac{n+n''}{n'}\mathpunct{\raisebox{.45ex}{,}}
\]
shows that, given the assumption that $n,n''$ are proportional to $\theta,\theta''$, the defect in the
method proposed in the preceding article arises only from the \textit{further} assumption
that $n'$ is proportional to $\theta'$. This replaces the factor $\dfrac{n+n''}{n'}$ with the less accurate value $\dfrac{\theta+\theta''}{\theta'}=1$,
from which the true value
\[
=1+\frac{\theta\theta''}{2\eta\eta''rr'r''\cos f\cos f'\cos f''}
\]
differs by a quantity of the second order (Article 128).
